\section{The shared open observer}
\label{sec:open}

By \eqref{eq:odd} the same run gives $\Open(2n-1)=\Closed(n)$. For $\Open(2n)$
we need, by \eqref{eq:even}, the sum of $e(Q)$ over the connected pairs. The
following identity localizes exterior arches to events visible in the inward
words.

\begin{lemma}[Exterior cut identity]
\label{lem:exterior}
\lean{Meanders.FirstCrossing.exteriorArchCutEquiv, Meanders.FirstCrossing.exteriorCount_cut, Meanders.FirstCrossing.exteriorCut_reflection}
\leanok
\uses{lem:cut}
Let $Q\in\NC_n$, let $0\le L\le2n$ and $w=h_Q(L)$. Let $r_L$ and $r_R$ be the
numbers of $D$ letters in the left and right inward words of $Q$ whose
\emph{original} pre-step height is $1$ (equivalently, whose post-step height is
$0$). Then
\[
 e(Q)=r_L+r_R+\ind{w>0}.
\]
Moreover the exterior arches are in bijection with these $r_L+r_R$ letters
together with, when $w>0$, one distinguished through arch.
\end{lemma}

\begin{proof}
\leanok
Every arch of $Q$ lies entirely in the left half, entirely in the right half,
or crosses the cut. An arch inside the left half is exterior if and only if no
arch encloses it. In the stack reading of the left word the arch closes at a
$D$ letter; it is enclosed by some arch exactly when another opening is still
on the stack at that moment, that is, when the pre-step height exceeds $1$
(an older opening on the stack either closes later, enclosing the arch, or is a
through opening, which also encloses it). The right half is symmetric, since
reflection preserves nesting. Through arches are pairwise nested
(\cref{lem:cut}); if $w>0$ the outermost one, joining the two oldest unmatched
openings, is enclosed by nothing, and every other through arch is enclosed by
it, while an internal arch cannot enclose an arch that crosses the cut. The
three classes are disjoint and exhaustive.
\end{proof}

\begin{definition}[Observer]
\label{def:observer}
\lean{Meanders.FirstCrossing.lowerIncrement, Meanders.FirstCrossing.Sector.lowerBonus}
\leanok
\uses{lem:partition, lem:exterior}
Each key carries a pair of integers $(C,E)$, seeded at $(1,0)$. When a point on
side $s$ is processed with a $D$ letter for $Q$ and original pre-step height
$h_{Q,s}=i_s-2c_{Qs}=1$, set $a=1$; otherwise $a=0$. The transition sends
$(C,E)$ to $(C,\,E+aC)$, merging equal keys by componentwise addition. At
acceptance a \HIGH\ sector emits $(C,\,E+bC)$ with $b=\ind{v>0}$ (which is
$1$ in every nonempty sector by \cref{lem:partition}); \LOW\ emits $(C,E)$.
\end{definition}

\begin{theorem}[Joint correctness]
\label{thm:joint}
\lean{Meanders.FirstCrossing.jointLayer_invariant, Meanders.FirstCrossing.AcceptedWord.event_exterior, Meanders.FirstCrossing.evaluateJoint_correct, Meanders.FirstCrossing.count_eq}
\leanok
\uses{lem:partition, lem:filter, lem:future, lem:cycle, lem:exterior, def:observer}
For $n\ge1$ and $1\le K\le n$, the sum over \LOW\ and all admissible \HIGH\
sectors of the emitted first components is $\Closed(n)$, and the sum of the
emitted second components is $\Open(2n)$.
\end{theorem}

\begin{proof}
\leanok
Fix a sector and a stage $t$. For a key $\kappa$ let $H_t(\kappa)$ be the set of
surviving partial source words with key $\kappa$, and for such a word $x$ let
$r(x)$ be the number of $Q$-letters $D$ read so far from original height $1$.
We claim that the stored pair equals
$\bigl(|H_t(\kappa)|,\ \sum_{x\in H_t(\kappa)}r(x)\bigr)$. This holds for the
seed. The increment $a$ of the next letter is a function of the key, the stage
and the label, because $h_{Q,s}=i_s-2c_{Qs}$ is read from the counters
(\cref{lem:filter}); extending each $x$ by a label with increment $a$ changes
$r$ to $r+a$, and the multiset of successors is the same for every $x$ in
$H_t(\kappa)$ by \cref{lem:future}. Summing gives $(C,E+aC)$ for each label,
and merging equal targets adds the sums. Rejected words contribute to neither
channel. At the last stage, an accepted word $x$ is a connected pair
(\cref{lem:cycle}) and $r(x)=r_L+r_R$, so $r(x)+b=e(Q)$ by
\cref{lem:exterior} with the sector's cut $L$ and $w=v$ (in \LOW\ the cut is
$L=2n$, $w=0$). Summing over the partition of \cref{lem:partition} and applying
\eqref{eq:even} completes the proof.
\end{proof}

\begin{corollary}
\label{cor:open}
\lean{Meanders.FirstCrossing.count_open_even, Meanders.FirstCrossing.count_open_odd}
\uses{thm:joint, thm:cost, cor:balance}
One run at order $n$ returns $\bigl(\Closed(n),\Open(2n-1),\Open(2n)\bigr)$
within the bounds of \cref{thm:cost,cor:balance}; the observer adds two
$O(n)$-bit additions per transition and no key field, since $0\le E\le nC$
(a surviving history contains at most $n$ lower $D$ letters, each
contributing at most $C$) and the terminal bonus adds at most $C$. Hence $\Open(m)$ is
computable in $\Ostar(2^{\lceil m/2\rceil})$ for every $m\ge1$, and
$\Open(0)=1$.
\end{corollary}

\begin{remark}[Two shortcuts that fail]
\label{rem:shortcuts}
(a) The event must use the \emph{original} height, not the number of retained
ports of the group. At $n=3$ in sector $(L,u,v)=(2,2,2)$ there is a state
before the last right point whose lower-right original height is $3$ while its
retained group has one port; a lower $D$ there contributes nothing, and a test
on retained length would wrongly add one.
(b) The lower channel cannot be replaced by half of a symmetric upper-plus-lower
channel sector by sector: at $n=8$ the sectors $(6,4,6)$ and $(6,6,4)$ both
contain $506$ connected pairs but have lower exterior totals $701$ and
$1021$. Exchanging $P$ and $Q$ maps a sector to a different sector, so the
symmetric total is $2\Open(2n)$ only after summing all sectors.
\end{remark}
